Let $X$ be a locally compact Hausdorff space and $H$ be a Hausdorff space. We denote by $C(X,H)$ the set of continuous functions from $X$ to $H$. We also fix for the remainder of this section a function $\theta\in C(X,H)$\nomenclature[Z]{$C(X,H)$}{Set of continuous functions}. For two functions $f,g:X\to H$, we denote\nomenclature[Z]{$[f\neq\theta],[f=\theta]$}{}
\[[f\neq g]=\left\{x\in X:f(x)\neq g(x)\right\}\quad\text{and}\quad[f=g]=X\setminus[f\neq g].\]

\begin{definition}\label{definitionsigma}\index{Support}\nomenclature[Z]{$\supp(f),\supp^\theta(f)$}{Support of $f$}
Given $f\in C(X,H)$, we define the ($\theta$-)\emph{support} of $f$ as
\[\supp^\theta(f)=\overline{[f\neq\theta]}.\]
We also denote by $\sigma^\theta(f)$ the interior of $\supp^\theta(f)$, and $Z^\theta(f)$ the complement of $\supp^\theta(f)$:\nomenclature[Z]{$\sigma(f),\sigma^\theta(f)$}{Interior of $\supp(f)$}\nomenclature[Z]{$Z(f),Z^\theta(f)$}{Complement of $\supp(f)$}
\[\sigma^\theta(f)=\operatorname{int}\supp^\theta(f)\qquad\text{and}\qquad Z^\theta(f)=X\setminus\supp^\theta(f).\]
Let $C_c(X,\theta)$ be the set of continuous functions from $X$ to $H$ with compact ($\theta$-)support. Whenever there is no risk of confusion, we will drop $\theta$ from the notation and write simply $\supp(f)$, $\sigma(f)$, $Z(f)$ and $C_c(X)$.\nomenclature[Z]{$C_c(X),C_c(X,\theta)$}{Set of compactly supported functions}

For any two functions $f,g\in C(X,H)$, we define
\begin{enumerate}
\item $f\perp g$: if $[f\neq\theta]\cap[g\neq\theta]=\varnothing$; we say that $f$ and $g$ are \emph{disjoint};\index{Disjoint functions}\nomenclature[Z]{$f\perp g$}{Disjointness}
\item $f\perpp g$: if $\supp(f)\cap\supp(g)=\varnothing$; we say that $f$ and $g$ are \emph{strongly disjoint};\index{Disjoint functions!Strongly disjoint}\nomenclature[Z]{$f\perpp g$}{Strong disjointness}
\item $f\subseteq g$: if $\sigma(f)\subseteq\sigma(g)$;\nomenclature[Z]{$f\subseteq g$}{}
\item $f\Subset g$: if $\supp(f)\subseteq\sigma(g)$.\nomenclature[Z]{$f\Subset g$}{}
\end{enumerate}
\end{definition}

Note that if $f\in C_c(X)$, then $f\Subset g$ if and only if $\sigma(f)\Subset\sigma(g)$, where $\Subset$ is the \emph{compact containment} relation (Definition \ref{definitioncompactcontainment})

\begin{remark}
In general, $\sigma(f)$ is the regularization\ftnt{The \emph{regularization} of an open set $A$ in a topological space $X$ is $\operatorname{int}\overline{A}$, and it is the smallest regular open set of $X$ containing $A$} of $[f\neq\theta]$, and so $\sigma(f)$ may ``forget'' some points on which $f(x)=\theta(x)$. For example, if $X=H=[0,1]$, $\theta=0$ and $f(x)=x$, then $\sigma(f)=X$ even though $f(0)=0$.

Also, $Z(f)$ is the complement of $\sigma(f)$ in the lattice of regular open sets of $X$ (see \cite[Chapter 10]{MR2466574}), and note that $\overline{\sigma(f)}=\supp(f)$.
\end{remark}

\begin{example}
If $H=\mathbb{R}$ or $\mathbb{C}$, and $\theta=0$ is the constant zero function, we obtain the usual notions of support, and $f\perp g$ if and only if $fg=0$.
\end{example}

\begin{example}
If $H=X$ and $\theta=\operatorname{id}_X$, then $\supp(f)=\overline{\left\{x\in X:f(x)\neq x\right\}}$.
\end{example}

We will be interested in relations between $\subseteq,\Subset,\perp$ and $\perpp$.

\begin{lemma}\label{lemmadisjointnessintermsofsigma}
$f\perp g$ if and only if $\sigma(f)\cap\sigma(g)=\varnothing$.
\end{lemma}
\begin{proof}
By definition, $f\perp g\iff[f\neq\theta]\cap[g\neq\theta]=\varnothing$. Since both $[f\neq\theta]$ and $[g\neq\theta]$ are open (because $H$ is Hausdorff), we can take closures and interiors of $[f\neq\theta]$ and obtain
\[f\perp g\iff\sigma(f)\cap[g\neq\theta]=\varnothing,\]
and repeating the argument with $[g\neq\theta]$ yields the result.\qedhere
\end{proof}

We will, moreover, be interested in recovering a locally compact Hausdorff topological space $X$ not from the whole set $C_c(X)$, but instead from a subcollection $\mathcal{A}\subseteq C_c(X)$. We will need to assume, however, that there are enough functions in $\mathcal{A}$ in order to separate points of $X$.

\begin{definition}\label{definitionregularperpp}
Let $\mathcal{A}\subseteq C_c(X)$ be a subset containing $\theta$. Denote $\sigma(\mathcal{A})=\left\{\sigma(f):f\in\mathcal{A}\right\}$. We say that $(X,\theta,\mathcal{A})$ (or simply $\mathcal{A}$) is
\begin{enumerate}
\item \emph{weakly regular} if $\sigma(\mathcal{A})$ is a basis for the topology of $X$.\index{Regular family of functions!weakly regular}
\item \emph{regular} if for every $x\in X$, every neighbourhood $U$ of $x$ and every $c\in H$ there is $f\in\mathcal{A}$ with $f(x)=c$ and $\supp(f)\subseteq U$.\index{Regular family of functions}
%\item \emph{normal} if for every compact $K\subseteq X$, every closed $F\subseteq X$, and every pair of elements $\alpha,\beta\in H$, there exists $f\in\mathcal{A}$ with $f=\alpha$ on $K$ and $\beta$ on $F$.
%\item \emph{directed} if there exists a function $S:\mathcal{A}^2\to\mathcal{A}$ with the property that
%\[[S(f,g)=\theta]=[f=\theta]\cap[g=\theta]\]
\end{enumerate}
\end{definition}

\begin{definition}\label{definitiondisjointcover}\index{Cover}
Suppose $\mathcal{A}\subseteq C_c(X)$ is weakly regular. A family $A\subseteq\mathcal{A}$ is a \emph{cover} of an element $b\in\mathcal{A}$ if given $h\in\mathcal{A}$, $h\perp a$ for all $a\in A$ implies $h\perp b$.
\end{definition}

By Lemma \ref{lemmadisjointnessintermsofsigma}, $A$ is a cover for $b$ if and only if $\left\{\sigma(f):f\in A\right\}$ is a cover for $\sigma(b)$, in the sense of Definition \ref{definitioncoverinposet}.

\begin{lemma}\label{lemmacovers}
Suppose $\mathcal{A}$ is weakly regular, and let $A\subseteq\mathcal{A}$ and $b\in\mathcal{A}$. The following are equivalent:
\begin{enumerate}
\item[(1)] $A$ is a cover of $b$;
\item[(2)] The closure of $\bigcup_{a\in A}[a\neq\theta]$ contains $\supp(b)$.
\end{enumerate}
\end{lemma}
\begin{proof}
(1)$\Rightarrow$(2): Let $x\in\supp(b)$. Take an open neighbourhood of $x$ of the form $\sigma(h)$, $h\in\mathcal{A}$. Since $\supp(b)=\overline{\sigma(b)}$, the intersection $\sigma(h)\cap\sigma(b)$ is nonempty and thus $h$ and $b$ are not disjoint. From $A$ being a cover, $h$ is not disjoint to some $a\in A$, which means that $\sigma(h)\cap[a\neq\theta]$ is nonempty. Since $A$ is weakly regular then $x$ is in the closure of $\bigcup_{a\in A}[a\neq\theta]$.

(2)$\Rightarrow$(1): Suppose $h\in\mathcal{A}$ is such that $h\perp a$ for all $a\in A$. This means that $(\bigcup_{a\in A}[a\neq\theta])\cap[h\neq\theta]=\varnothing$. Taking the closure of the first term and using (2) we conclude that $[b\neq\theta]\cap[h\neq\theta]\subseteq\supp(b)\cap[h\neq\theta]=\varnothing$, so $h\perp b$.\qedhere
\end{proof}

%\begin{lemma}\label{lemmafunctionsoncompacts}
%Suppose $\mathcal{A}$ is weakly regular and directed, $K\subseteq X$ is compact, $F\subseteq X$ is closed and $K\cap F=\varnothing$. Then there exists $h\in\mathcal{A}$ such that $K\subseteq\sigma(h)$ and $\supph\cap F=\varnothing$.
%\end{lemma}
%\begin{proof}
%From regularity and compactness, there are $h_1,\ldots,h_n$ such that $K\subseteq\bigcup_{i=1}^n\sigma(h_i)$ and $\supp(h_i)\cap F=\varnothing$. Letting $h=S(h_1,S(h_2,\cdots))$ we are done.\qedhere
%\end{proof}

If $\mathcal{A}\subseteq C_c(X)$ and $\theta\in\mathcal{A}$, note that $\subseteq$ is a preorder on $\mathcal{A}$, and the only infimum\ftnt{Infima for preorders are defined the same way as for orders, namely if $A$ is a subset of a preordered space $(P,\leq)$, an element $x$ is called an infimum of $A$ if $x\leq a$ for all $a\in A$, and if for any other element $y$ such that $y\leq a$ for all $a\in A$, we have $y\leq x$.} of $\mathcal{A}$ is $\theta$. Alternatively, $\theta$ is the only element of $\mathcal{A}$ such that $\theta\perp\theta$. This shows that the function $\theta$ is uniquely determined in terms of either $\perp$ or $\subseteq$. 

\begin{theorem}\label{theoremrelationsrelations}
Suppose $\mathcal{A}$ is weakly regular. If $f,g\in\mathcal{A}$, then
\begin{enumerate}
\item[(a)] $f\subseteq g\iff\forall h(h\perp g\Rightarrow h\perp f)$;
\item[(b)] $f\perp g\iff$ The $\subseteq$-infimum of $\left\{f,g\right\}$ is $\theta$ (which is the $\subseteq$-infimum of $\mathcal{A}$);
\item[(c)] $f\subseteq g\iff\forall h(h\Subset f\Rightarrow h\Subset g)$;
\item[(d)] $f\subseteq g\iff\forall h(h\perpp g\Rightarrow h\perpp f)$;
\item[(e)] $f\perpp g\iff\exists h_1,k_1,\ldots,h_n,k_n\in\mathcal{A}$ such that $\left\{h_1,\ldots,h_n\right\}$ is a cover of $f$, $h_i\Subset k_i$ and $k_i\perp g$;
\item[(f)] $f\Subset g\iff\forall b\in\mathcal{A}$, $\exists h_1,\ldots,h_n\in\mathcal{A}$ such that $\left\{h_1,\ldots,h_n,g\right\}$ is a cover of $b$ and $h_i\perpp f$.
\end{enumerate}
By items (a) and (b), $\perp$ and $\subseteq$ are equi-expressible (i.e., each one is completely determined by the other). By (c) and (d) one can recover $\subseteq$ (and hence $\perp$) from either $\Subset$ or $\perpp$, which in turn implies, from (e) and (f), that $\Subset$ and $\perpp$ are also equi-expressible.
\end{theorem}
\begin{proof}
\begin{enumerate}
\item[(a)] $\Rightarrow$: If $f\subseteq g$ and $h\perp g$, then $\sigma(f)\cap\sigma(h)\subseteq\sigma(g)\cap\sigma(h)=\varnothing$, thus $h\perp f$.

$\Leftarrow$: Suppose $\sigma(f)$ is not contained in $\sigma(g)$. Since $\sigma(f)=\operatorname{int}([f\neq g])$ and $\sigma(g)$ is regular, this means that $[f\neq\theta]$ is not contained in $\overline{\sigma(g)}=\supp(g)$, so there is an open  set $U$ such that $f\neq\theta$ on $U$ but $g=\theta$ on $U$. By weak regularity, there is some $h$ such that $\varnothing\neq[h\neq\theta]\subseteq U$, so $f$ and $h$ are not disjoint but $g$ and $h$ are.

\item[(b)] $\Rightarrow$: If $h\subseteq f,g$, then $\sigma(h)\subseteq\sigma(f)\cap\sigma(g)=\varnothing$, so $h=\theta$.

$\Leftarrow$: If $f$ and $g$ are not disjoint, then there is some $h$ with $\varnothing\neq\sigma(h)\subseteq\sigma(f)\cap\sigma(g)$, by weak regularity, so $h\subseteq f,g$ and $h\neq\theta$.

\item[(c)] This is immediate since $X$ is locally compact Hausdorff and $\mathcal{A}$ is weakly regular, so $\sigma(f)=\bigcup\left\{\sigma(h):h\Subset f\right\}$ for all $f$.

\item[(d)] The proof is essentially the same as (a), so we omit it. %$\Rightarrow$: If $f\subseteq g$ and $h\perpp g$, then for $x\in\supp(h)$ there is a neighbourhood of $x$ on which $g=c$, and so $f=c$ as well, which implies $h\perpp f$.

%$\Leftarrow$: Same as in (i), but just make sure to take $h$ with $\supph\subseteq\mathcal{A}$.

\item[(e)] $\Rightarrow$: Suppose $f\perpp g$. For every $x\in\supp(f)$, take $h_x,k_x\in\mathcal{A}$ such that $x\in\sigma(h_x)$, $h_x\Subset k_x$ and $k_x\perp g$. By compactness, we can take finitely many points $x_1,\ldots,x_n$, and associated functions $h_1,k_1,\ldots,h_n,k_n$, such that $\left\{h_1,\ldots,h_n\right\}$ covers $f$.

$\Leftarrow$: Suppose such $h_i,k_i$ exist. Then
\[\supp(f)\subseteq\bigcup_{i=1}^n\supp(h_i)\subseteq\bigcup_{i=1}^n\sigma(k_i)\subseteq X\setminus\supp(g)\]
and so $f\perpp g$.
\item[(f)] $\Rightarrow$: Suppose $f\Subset g$ and take any $b\in\mathcal{A}$. Since $\supp(b)\setminus\sigma(g)$ is compact and does not intersect $\supp(f)$, we can take $h_1,\ldots,h_n\in\mathcal{A}$ such that $h_i\perpp f$ and $\supp(b)\setminus\sigma(g)\subseteq\bigcup_i\sigma(h_i)$, which implies that $\left\{h_1,\ldots,h_n,g\right\}$ is a cover of $b$.

$\Leftarrow$: Suppose the second condition holds. Again by compactness, take $b_1,\ldots,b_M$ in $\mathcal{A}$ such that $\supp(f)\cup\supp(g)\subseteq\bigcup_{k=1}^M\sigma(b_k)$. For each $k$ take functions $h_i^k$ as in the given condition for $b_k$.

Given $k$, we have $\sigma(b_k)\subseteq\bigcup_i\supp(h_i^k)\cup\supp(g)$, so taking complements,
\[\bigcap_iZ(h_i^k)\cap Z(g)\cap\sigma(b_k)=\varnothing,\]
or equivalently $\bigcap_iZ(h_i^k)\cap\sigma(b_k)\subseteq\sigma(b_k)\setminus Z(g)\subseteq\supp(g)$. Taking interiors on both sides yields $\bigcap_iZ(h_i^k)\cap\sigma(b_k)\subseteq\sigma(g)$.

Now from $h_i^j\perpp f$ we obtain
\begin{align*}
\supp(f)&\subseteq\bigcap_{i,j}Z(h_i^j)\cap\bigcup_{k=1}^M\sigma(b_k)=\bigcup_{k=1}^M\left[\bigcap_{i,j}Z(h_i^j)\cap\sigma(b_k)\right]\\
&\subseteq\bigcup_{k=1}^M\left[\bigcap_i Z(h_i^k)\cap\sigma(b_k)\right]\subseteq \sigma(g),
\end{align*}
so $f\Subset g$.\qedhere
\end{enumerate}
\end{proof}

\begin{remark}
One should be careful with the connections between the pairs of relations $(\perp,\perpp)$ and $(\subseteq,\Subset)$. For example, $\perp$ and $\perpp$ can coincide but $\subseteq$ and $\Subset$ may not and vice-versa. See the example below.
\end{remark}

\begin{example}
Let $X=H=\mathbb{R}$ and $\theta=0$, so that we are dealing with the usual notion of support. Let $\left\{(a_n,b_n):n\in\mathbb{N}\right\}$ ($a_n<b_n$) be a basis of open intervals for the usual topology of $\mathbb{R}$ with $|b_n-a_n|\to 0$. Let $\left\{p_n:n\in\mathbb{N}\right\}$ be an one-to-one enumeration of the prime numbers.

For each $n$, let $\widetilde{a_n}$ and $\widetilde{b_n}$ be, respectively, the largest and smallest rational numbers with denominators $p_n$ as reduced fractions, and which satisfy
\[\widetilde{a}_n\leq a_n<b_n\leq \widetilde{b}_n.\]
Namely,
\[\widetilde{a}_n=\frac{\lfloor a_np_n\rfloor}{p_n}\qquad\text{and}\qquad\widetilde{b}_n=\frac{\lceil b_np_n\rceil}{p_n}\]
where $\lfloor\cdot\rfloor$ and $\lceil\cdot\rceil$ are the usual \emph{floor} and \emph{ceiling} functions. In particular, $|\widetilde{a_n}-a_n|+|\widetilde{b_n}-b_n|\leq 2/p_n\to 0$ (since the enumeration of the primes is one-to-one), and thus the sets $U_n=(\widetilde{a_n},\widetilde{b_n})$ also form a basis of $\mathbb{R}$.

For each $n$, let $f_n\in C_c(\mathbb{R})$ with $\sigma(f_n)=U_n$, e.g.
\[f_n(x)=\max(0,(x-\widetilde{a_n})(\widetilde{b_n}-x)),\]
and let $\mathcal{A}=\left\{f_n:n\in\mathbb{N}\right\}$. Then $\perp$ and $\perpp$ coincide on $\mathcal{A}$, as do $\subseteq$ and $\Subset$, since the boundaries of all $U_n$ are pairwise disjoint.

Letting $V=(\widetilde{a_1},\widetilde{b_1}+1)$ and $g_V$ be a continuous function with $\sigma(g_V)=V$, then $\perp$ and $\perpp$ still coincide in $\mathcal{A}\cup\left\{g_V\right\}$, however $\subseteq$ and $\Subset$ do not, since $f_1\subseteq g_V$ but $f_1\not\Subset g_V$.

Alternatively, set $W=(\widetilde{b_1},\widetilde{b_1}+1)$ and let $g_W$ be any continuous function with $\sigma(g_W)=W$. Then $\subseteq$ and $\Subset$ still coincide in $\mathcal{A}\cup\left\{g_W\right\}$, however $\perp$ and $\perpp$ do not, because $f_1\perp g$ but not $f_1\perpp g$.
\end{example}

\begin{example}[{Kania-Rmoutil, \cite{MR3813611}}]\label{examplekaniarmoutil}
Suppose that $X$ is a Stone space (Definition \ref{definitionstonespace}) and that $H$ is a Hausdorff space with at least two points. Let $\mathcal{A}$ consist of all continuous $H$-valued functions with finite range and $\theta$ be any element of $\mathcal{A}$. Regularity of $\mathcal{A}$ is an immediate consequences of $X$ being zero-dimensional (and $H$ containing more than one point). Define the \emph{compatibility ordering} on $C(X)$ by
\[f\preceq g\iff g|_{\supp^\theta(f)}=f|_{\supp^\theta(f)}.\]
Then $\perp$ coincides with $\perpp$, and it can be described in terms of $\preceq$ by
\[f\perp g\iff f\perpp g\iff\inf_{\preceq}(f,g)=\theta(=\inf_{\preceq}\mathcal{A})\text{ and }\left\{f,g\right\}\text{ has a }\preceq\text{-upper bound.}\]
\end{example}

%\begin{example}\label{exampleliwong}
%If $H=\mathbb{R}$ or $\mathbb{C}$ and $X$ is regular, then $(C(X,H\setminus 0),1)$ is regular and directed via $S(f,g)=|f-1|+|g-1|+1$.
%\end{example}

%\section{The main theorems}

We now state and prove the two main theorems of this chapter. Fix two locally compact Hausdorff spaces $X$ and $Y$, and for $Z\in\left\{X,Y\right\}$ a Hausdorff space $H_Z$, a continuous map $\theta_Z:Z\to H_Z$, and a subset $\mathcal{A}(Z)\subseteq C_c(Z,\theta_Z)$.

\begin{definition}
We call a map $T:\mathcal{A}(X)\to\mathcal{A}(Y)$ a \emph{$\perpp$-morphism} if $f\perpp g$ implies $Tf\perpp Tg$; $T$ is a \emph{$\perpp$-isomorphism} if it is bijective and both $T$ and $T^{-1}$ are $\perpp$-morphisms. $\perp$, $\subseteq$ and $\Subset$-isomorphisms are define analogously.\index{$\perpp$-morphism}\index{$\perp$-morphism}\index{$\subseteq$-morphism}\index{$\Subset$-morphism}
\end{definition}

%\begin{lemma}\label{lemmasemilattice}
%$\sigma(A)$ is a $\lor$-semilattice, or more precisely a sub-$\lor$-semilattice of regular open sets\footnote{The regular open sets of a compact Hausdorff space form a lattice, with join operation given by $A\lor B=\operatorname{int}(\overline{A\cup B})$.}.
%\end{lemma}
%\begin{proof}
%Consider a direction map $S:A^2\to A$. We need to show that $\sigma(S(f,g))=\operatorname{int}(\overline{\sigma(f)\cup\sigma(g)})$. We have
%\[[S(f,g)\neq c]=[f\neq c]\cup[g\neq c]\]
%so taking closures and interiors twice yields $\operatorname{int}(\overline{\sigma(f)\cup\sigma(g)})\subseteq\sigma(S(f,g))$. Conversely,
%\[[S(f,g)\neq c]=[f\neq c]\cup[g\neq c]\subseteq \sigma(f)\cup\sigma(g)\subseteq\operatorname{int}(\overline{\sigma(f)\cup\sigma(g)}),\]
%and taking closures and interiors again yields the converse inclusion.\qedhere
%\end{proof}

%\begin{theorem}\label{welldefinedorderisomorphism}
%Suppose $T:A_X\to A_Y$ is a $\perp$-isomorphism. Then for all $f\in A_X$, $\sigma(f)\subseteq\sigma(g)$ if and only if $\sigma(Tf)\subseteq\sigma(Tg)$. In particular, $\sigma(A_X)$ and $\sigma(A_Y)$ are isomorphic.
%\end{theorem}
%\begin{proof}
%We are done if we describe $\sigma(f)\subseteq\sigma(g)$ in terms of $\perp$\footnote{No, because $\sigma(A)$ is not closed under union, just closed under regular open joins...}. First note that $\sigma(f)\subseteq\sigma(g)$ if and only if $\sigma(f)\subseteq\suppg$. From this fact and regularity, we obtain
%\[f\perp g\iff\forall h\in A, [h\neq c]\cap\suppg=\varnothing\iff [h\neq c]\cap\sigma(f)=\varnothing\]
%Since $\suppg=\overline{[g\neq c]}$ and $\suppf=\overline{[f\neq c]}$ is regular, we can rewrite this as
%\begin{align*}
%f\perp g&\iff\forall h\in A, [h\neq c]\cap[g\neq c]\Rightarrow[h\neq c]\cap[f\neq c]=\varnothing\\
%&\iff \forall h\in A, h\perp g\Rightarrow h\perp f.\qedhere
%\end{align*}
%\end{proof}

\begin{theorem}\label{theoremdisjoint}
Suppose $\mathcal{A}(X)$ and $\mathcal{A}(Y)$ are weakly regular and $T:\mathcal{A}(X)\to\mathcal{A}(Y)$ is a $\perp$-isomorphism. Let $f,g\in\mathcal{A}(X)$. Then $\sigma(f)\subseteq \sigma(g)$ if and only if $\sigma(Tf)\subseteq \sigma(Tg)$. In particular, $Z(f)=\varnothing$ if and only if $Z(Tf)=\varnothing$.
\end{theorem}
\begin{proof}
Immediate from \ref{theoremrelationsrelations}(a), because every $\perp$-morphism is a $\subseteq$-morphism (and vice-versa).\qedhere
\end{proof}

\begin{theorem}\label{maintheorem}
If $\mathcal{A}(X)$ and $\mathcal{A}(Y)$ are weakly regular and $T:\mathcal{A}(X)\to\mathcal{A}(Y)$ is a $\perpp$-isomorphism then there is a unique homeomorphism $\phi:Y\to X$ such that $\phi(\supp(Tf))=\supp(f)$ for all $f\in \mathcal{A}(X)$ (equivalently, $\phi(\sigma(Tf))=\sigma(f)$, or $\phi(Z(Tf))=Z(f)$, for all $f\in\mathcal{A}(X)$).
\end{theorem}
\begin{proof}
By Theorem \ref{theoremrelationsrelations}, $\perpp$-isomorphisms coincide with $\Subset$-isomorphisms, so the result is an easy application of Theorem \ref{theoremrecoverxfromultrafilters}. We add the details for the sake of completeness.

Let $\Phi:\sigma(\mathcal{A}(X))\to\sigma(\mathcal{A}(Y))$ be the map given by $\Phi(\sigma(f))=\sigma(Tf)$ for all $f\in\mathcal{A}(X)$, which is well-defined by \ref{theoremdisjoint}. This is a $\Subset$-isomorphism, so, using the same notation as in \ref{theoremrecoverxfromultrafilters}, we induce a homeomorphism
\[\Phi^*:\widehat{\sigma(\mathcal{A}(Y))}^{\Subset}\to\widehat{\sigma(\mathcal{A}(X))}^{\Subset},\quad\Phi^*(F)=\Phi^{-1}(F).\]
For $Z\in\left\{X,Y\right\}$, let $\kappa_Z:Z\to\widehat{\sigma(\mathcal{A}(Z))}^{\Subset}$ be the homeomorphism given by
$\kappa_Z(z)=\left\{\sigma(f):z\in\sigma(f)\right\}$. We then define $\phi=\kappa_X^{-1}\circ\Phi^*\circ\kappa_Y:Y\to X$, which is a homeomorphism. For all $y\in Y$,
\[\kappa_X(\phi(y))=\Phi^{-1}(\kappa_Y(y))=\left\{\sigma(f):y\in\sigma(Tf)\right\}\]
which means that for each $f\in\mathcal{A}(X)$, $y\in\sigma(Tf)$ if and only if $\phi(y)\in\sigma(f)$. Therefore $\phi(\sigma(Tf))=\sigma(f)$, or equivalently taking closures we obtain $\phi(\supp(Tf))=\supp(f)$.\qedhere
\end{proof}

For an alternative proof, see the end of Section \ref{sectionperppideals}.

\begin{definition}\label{definitionthomeomorphism}\index{$T$-homeomorphism}
The unique homeomorphism $\phi$ associated with $T$ as in \ref{maintheorem} will be called the \emph{$T$-homeomorphism}.
\end{definition}

We finish this section by proving that the hypothesis that $T$ is a $\perpp$-isomorphism in Theorem \ref{maintheorem} cannot be weakened to a $\perp$-isomorphism in general (Corollaries \ref{corollary01stoneperp} and \ref{corollary01circleperp}).

Let $X$ and $Y$ be locally compact Hausdorff spaces. We will consider only real-valued functions and $\theta_X$ and $\theta_Y$ the respective zero functions on $X$ and $Y$, so that supports of continuous functions are the usual ones. By Theorem \ref{theoremdisjoint}, any $\perp$-isomorphism $T:C_c(X)\to C_c(Y)$ induces an order isomorphism between collections of regular open sets of $X$ and $Y$, respectively. So in order to find $\perp$-isomorphisms, we will look at spaces with isomorphic collections of regular open sets.

\begin{definition}
	Let $X$ be a topological space. We denote by $\operatorname{RO}(X)$ the collection of regular open subsets of $X$, and by $\operatorname{RO}_K(X)$ the collection of regular open subsets of $X$ with compact closure. Given $A\in \operatorname{RO}_K(X)$, we define $\Sigma_X(A)=\left\{f\in C_c(X):\sigma(f)=A\right\}$.
\end{definition}

A generalized Boolean algebra is \emph{conditionally complete} if every bounded family admits a join. In the case of unital Boolean algebras we just call these \emph{complete} (see \cite[Definition I-4.5]{MR1975381}).

It is well-known that the family of regular open subsets of a topological space is a complete unital Boolean algebra (see \cite[Theorem 10.1]{MR2466574}). Similarly, one can prove that $\operatorname{RO}_K(X)$ is a conditionally complete generalized Boolean algebra for any Hausdorff space $X$.
	
In order to find non-homeomorphic spaces $X$ and $Y$ such that $C_c(X)$ and $C_c(Y)$ are $\perp$-isomorphic, we need the following lemma:

\begin{lemma}\label{lemmacardinalityofbigsigma}
	Let $X$ be a separable, locally compact Hausdorff space and $A\in\operatorname{RO}_K(x)$. If $A\neq\varnothing$ and the set $\Sigma_X(A)=\left\{f\in C_c(X):\sigma(f)=A\right\}$ is nonempty, then it has the cardinality of the continuum, $|\Sigma_X(A)|=2^{\aleph_0}$.
\end{lemma}
\begin{proof}
	If $f\in\Sigma_X(A)$ then $f\neq 0$, and the continuum-many functions $\lambda f$, $\lambda\in\mathbb{R}$, also belong to $\Sigma_X(A)$, so $|\Sigma_X(A)|\geq 2^{\aleph_0}$. Conversely, if $D\subseteq X$ is a dense countable subset of $X$ then the map $C(X)\to\mathbb{R}^D$, $f\mapsto f|_D$ is injective, so
	\[|\Sigma_X(A)|\leq|C(X)|\leq|\mathbb{R}^D|=|\mathbb{R}|^{|D|}=\left(2^{\aleph_0}\right)^{\aleph_0}=2^{\aleph_0\times\aleph_0}=2^{\aleph_0}.\qedhere\]
\end{proof}

\begin{theorem}\label{theoremperpisnotenough}
	Suppose that:
	\begin{enumerate}
		\item $X$ and $Y$ are separable, locally compact Hausdorff spaces;
		\item For all nonempty $A\in\operatorname{RO}_K(X)$ and $B\in\operatorname{RO}_K(Y)$, $\Sigma_X(A)$ and $\Sigma_Y(B)$ are nonempty;
		\item $\varphi:\operatorname{RO}_K(X)\to\operatorname{RO}_K(Y)$ is an order isomorphism (with respect to set inclusion).
	\end{enumerate}
	Then $C_c(X)$ and $C_c(Y)$ are $\perp$-isomorphic.
\end{theorem}
\begin{proof}
	Given $A\in\operatorname{RO}_K(X)$, the sets  $\Sigma_X(A)$ and $\Sigma_Y(\varphi(A))$ have the same cardinality, by Lemma \ref{lemmacardinalityofbigsigma}, so consider any bijection $T_A:\Sigma_X(A)\to\Sigma_Y(\varphi(A))$. Then the map
	\[T:C_c(X)\to C_c(Y),\qquad T(f)=T_{\sigma(f)}(f)\]
	is a $\perp$-isomorphism.\qedhere
\end{proof}

For our first example we employ the theory of Stone spaces and Stone duality as in Section \ref{sectionstoneduality}. Recall that
\begin{itemize}
	\item Given a zero-dimensional, locally compact Hausdorff space $\mathfrak{C}$, $\KB(\mathfrak{C})$ is the collection of compact-open subsets of $\mathfrak{C}$.
	\item Given a generalized Boolean algebra $B$, $\GP(B)$ is the topological space (unit groupoid) of ultrafilters on $B$, with the topology generated by the sets $[a]=\left\{F\in\GP(B):a\in F\right\}$, $a\in B$.
\end{itemize}

\begin{example}
	In general, arbitrary joins do not correspond to unions under Stone duality. Let $X$ be an infinite set, and consider the complete unital Boolean algebra $\mathcal{P}(X)$ of subsets of $X$, and $\mathfrak{C}=\GP(\mathcal{P}(X))$. By Proposition \ref{propositionnaturaltransformationsncsd}, the map $\mathcal{P}(X)\ni A\mapsto[A]\in\KB(\mathfrak{C})$ is an order isomorphism so $\KB(\mathfrak{C})$ is complete.
	
	Given $a\in X$, let us denote by $\left\{a\right\}^\uparrow=\left\{A\subseteq X:a\in A\right\}$ the (principal) ultrafilter generated by the element $\left\{a\right\}$ of $\mathcal{P}(X)$.
	
	Let $F\in\mathfrak{C}$ be any ultrafilter which does not contain finite subsets of $X$, and whose existence is guaranteed by Zorn's Lemma. Given a basic neighbourhood $[A]$ of $F$ in $\mathfrak{C}$, where $A\subseteq X$, choose any $a\in A$ and note that $\left\{a\right\}^\uparrow\in [A]\setminus\left\{F\right\}$, and therefore $F$ is non-isolated in $\mathfrak{C}$.
	
	Since $\mathfrak{C}$ is Hausdorff then $\bigcup\left\{[A]:F\not\in[A]\right\}=\mathfrak{C}\setminus\left\{F\right\}$, and as $F$ is non-isolated then this set is not clopen. Therefore the join of $\left\{[A]:F\not\in[A]\right\}$ in $\KB(\mathfrak{C})$ is actually $\mathfrak{C}$.
	
	More explicitly, in $\mathcal{P}(X)$ we have
	\[\bigvee\left\{A:F\not\in[A]\right\}=\bigvee\left\{A:A\not\in F\right\}\supseteq\bigvee_{x\in X}\left\{x\right\}=X,\]
	and thus $\bigvee\left\{[A]:F\not\in[A]\right\}=[X]=\mathfrak{C}$ in $\KB(\mathfrak{C})$.
\end{example}

\begin{lemma}[{Compare with \cite[Lemma 6.10]{MR1351251}}]\label{lemmawhenregularcompactopenareclopen}
Let $\mathfrak{C}$ be a zero-dimensional, locally compact Hausdorff space and suppose that the Boolean algebra $\KB(\mathfrak{C})$ is conditionally complete. Then $\KB(\mathfrak{C})=\operatorname{RO}_K(\mathfrak{C})$.
\end{lemma}
\begin{proof}
The inclusion $\KB(\mathfrak{C})\subseteq\operatorname{RO}_K(\mathfrak{C})$ is immediate, so let us deal with the converse.

Let $A\in\operatorname{RO}_K(\mathfrak{C})$, so in particular $A$ is contained in some element of $\KB(\mathfrak{C})$ and the family $\left\{V\in\KB(\mathfrak{C}):V\subseteq A\right\}$ is bounded. Let
\[U=\bigvee\left\{V\in\KB(\mathfrak{C}):V\subseteq A\right\}.\]
where the join is taken in $\KB(\mathfrak{C})$. We have
\[A=\bigcup\left\{V\in\KB(A):V\subseteq A\right\}\subseteq U.\]
If $W\in\KB(\mathfrak{C})$ and $W\subseteq U\setminus\overline{A}$, then $A\subseteq  U\setminus W$ and so
\[U=\bigvee\left\{V\in\KB(\mathfrak{C}):V\subseteq A\right\}\subseteq U\setminus W\]
hence $W=\varnothing$. This proves that $U\setminus\overline{A}=\varnothing$ and so $U\subseteq\overline{A}$. However, $U$ is clopen and $A$ is regular open, which imply $A=U\in\KB(\mathfrak{C})$.\qedhere
\end{proof}

\begin{lemma}\label{gprokseparableisseparable}
If $X$ is a separable locally compact Hausdorff space, then $\mathfrak{C}=\GP(\operatorname{RO}_K(X))$ is separable.
\end{lemma}
\begin{proof}
Let $\left\{x_n:n\in\mathbb{N}\right\}$ be a countable dense subset of $X$. For each $n$, by Zorn's Lemma there exists $G_n\in\mathfrak{C}$ such that
\[\left\{U\in\operatorname{RO}_K(X):x_n\in U\right\}\subseteq G_n.\]
Let us prove that $\left\{G_n:n\in\mathbb{N}\right\}$ is dense in $\mathfrak{C}$. Given a basic open subset $[A]$ of $\mathfrak{C}$, where $A\in\operatorname{RO}_K(X)$, find $n\in\mathbb{N}$ with $x_n\in A$, so $A\in G_n$ and therefore $G_n\in[A]$.\qedhere
\end{proof}

\begin{lemma}\label{lemmaregularopeninsecondcountableissigma}
If $X$ is a second-countable locally compact Hausdorff space and $A\in\operatorname{RO}_K(X)$, then there is $f\in C_c(X)$ such that $\sigma(f)=A$.
\end{lemma}
\begin{proof}
First choose a countable family of compact subsets $K_n\subseteq A$ such that $\bigcup_nK_n=A$. For each $n$ we can, by Urysohn's Lemma and regularity of $X$, find a continuous function $f_n:X\to[0,1]$ such that $f_n(k)=1$ for all $k\in K_n$ and $\supp(f_n)\subseteq A$. Letting $f=\sum_{n=1}^\infty 2^{-n}f_n$ we obtain $[f\neq 0]=\sigma(f)=A$, because $A$ is regular.\qedhere
\end{proof}

As a consequence of Lemmas \ref{lemmawhenregularcompactopenareclopen}, \ref{gprokseparableisseparable} and \ref{lemmaregularopeninsecondcountableissigma}, and Theorem \ref{theoremperpisnotenough}, we conclude:

\begin{corollary}\label{corollary01stoneperp}
	There exists a  zero-dimensional, compact Hausdorff topological space $\mathfrak{C}$ (name\-ly, $\mathfrak{C}=\GP(\operatorname{RO}_K([0,1]))$) -- which is, in particular, not homeomorphic to $[0,1]$ -- such that $C(\mathfrak{C})$ and $C([0,1])$ are $\perp$-isomorphic.
\end{corollary}

Note that, in the corollary above, the Boolean algebra $\operatorname{RO}([0,1])$ is uncountable, thus $\mathfrak{C}=\GP(\operatorname{RO}_K([0,1]))$ is not second-countable.

For our second example, we will consider only second-countable spaces. In the next lemma, we denote by $\operatorname{int}_Z$ and $\operatorname{cl}_Z$ the interior and closure operators on subsets of a topological space $Z$.

\begin{lemma}\label{lemmamorphismregularopenofsubset}
	Let $X$ be a topological space and $U$ an open set of $X$. Then
	\begin{enumerate}
		\item[(a)] If $A\in\operatorname{RO}(X)$ then $A\cap U\in\operatorname{RO}(U)$.
		\item[(b)] The map \[\varphi_U:\operatorname{RO}(X)\to\operatorname{RO}(U),\qquad \varphi_U(A)=A\cap U\]
		is order-preserving and surjective; The map $\zeta_U:A\mapsto \operatorname{int}_X(\operatorname{cl}_X(A))$ is an order-preserving right inverse to $\varphi_U$;
		\item[(c)] $\varphi_U$ is an order isomorphism if and only if $U$ is dense in $X$.
	\end{enumerate}
\end{lemma}
\begin{proof}
	\begin{enumerate}
		\item[(a)] Given $A\in\operatorname{RO}(X)$, since $U$ is open, we have
		\[\operatorname{int}_U(\operatorname{cl}_U(A\cap U))=\operatorname{int}_U(\operatorname{cl}_X(A)\cap U)=\operatorname{int}_X(\operatorname{cl}_X(A))\cap U=A\cap U,\]
		and this proves that $A\cap U\in\operatorname{RO}(U)$.
		\item[(b)] The last statement is the only non-trivial part. If $A\in\operatorname{RO}(U)$, we again use the fact that $U$ is open, as in item (a), to obtain
		\[\varphi_U(\zeta_U(A)))=\operatorname{int}_X(\operatorname{cl}_X(A))\cap U=\operatorname{int}_U(\operatorname{cl}_X(A)\cap U)=\operatorname{int}_U(\operatorname{cl}_U(A))=A,\]
		as desired.
		\item[(c)] If $U$ is not dense in $X$, then there exists a nonempty set $A\in\operatorname{RO}(X)$ which is disjoint with $U$, so
		\[\varphi_U(A)=\varnothing=\varphi_U(\varnothing),\]
		and thus $\varphi_U$ is not injective.
		
		Now assume that $U$ is dense in $X$, so let us prove that the map $\zeta_U$ of item (b) is a left inverse of $\varphi_U$. Given $A\in\operatorname{RO}(X)$, since $U$ is dense in $X$,
		\[\zeta_U(\varphi_U(A))=\operatorname{int}_X(\operatorname{cl}_X(A\cap U))=\operatorname{int}_X(\operatorname{cl}_X(A))=A.\qedhere\]
	\end{enumerate}
\end{proof}

Let $\mathbb{S}^1=\left\{z\in\mathbb{C}:|z|=1\right\}$ be the complex unit circle.

\begin{corollary}\label{corollary01circleperp}
	$C([0,1])$ and $C(\mathbb{S}^1)$ are $\perp$-isomorphic.
\end{corollary}
\begin{proof}
	Let $X=(0,1)$ and $Y=\mathbb{S}^1\setminus\left\{1\right\}$. Then $X$ and $Y$ are homeomorphic, and two applications of Lemma \ref{lemmamorphismregularopenofsubset} imply that $\operatorname{RO}([0,1])$ and $\operatorname{RO}(\mathbb{S}^1)$ are order-isomorphic. Lemmas \ref{gprokseparableisseparable} and \ref{lemmaregularopeninsecondcountableissigma}, and Theorem \ref{theoremperpisnotenough}, imply that $C([0,1])$ and $C(\mathbb{S}^1)$ are $\perp$-isomorphic.
\end{proof}
